% !TeX program = xelatex
% Overleaf: Menu/Settings > Compiler > XeLaTeX. The editor hint above does not
% override a project explicitly configured to use pdfLaTeX.
\documentclass[runningheads,a4paper]{llncs}
\usepackage[a4paper,textwidth=122mm,textheight=193mm,centering]{geometry}
\usepackage{fontspec}
\setmainfont{texgyretermes}[Extension=.otf,UprightFont=*-regular,BoldFont=*-bold,ItalicFont=*-italic,BoldItalicFont=*-bolditalic]
\setsansfont{texgyreheros}[Extension=.otf,UprightFont=*-regular,BoldFont=*-bold,ItalicFont=*-italic,BoldItalicFont=*-bolditalic]
\setmonofont{lmmono10-regular.otf}
\usepackage{amsmath,booktabs,tabularx,array}
\usepackage[unicode,hidelinks]{hyperref}
\urlstyle{same}
\setlength{\emergencystretch}{2em}
\renewcommand{\abstractname}{Abstract.}
\renewcommand{\keywordname}{\textbf{Keywords:}}
\renewcommand{\andname}{and}
\renewcommand{\lastandname}{, and}
\begin{document}
\title{Security Regression Testing of Indirect Prompt Injection Defenses under Versioned Agent Tool Contracts}
\titlerunning{Security Regression Testing for Tool-Using Agents}
\author{Nguyễn Thùy Dương \and Đặng Ngọc Huy \and Vũ Hồng Quang}
\authorrunning{Group 11}
\institute{Group 11\\Cybersecurity Governance Policy Course}
\maketitle
\begin{abstract}
LLM agents can invoke tools to read and modify data. Documents or tool outputs may contain untrusted instructions that redirect an agent beyond the user's task. This report presents a security regression-testing design for a tool contract migrating from version $v0$ to $v1$. The design holds the task, initial state, and agent configuration constant while measuring security and task utility as separate outcomes. The group implements task-scoped file authorization under frozen and explicitly migrated policy states and checks the environment before and after tool calls. Deterministic source checks show that the frozen policy rejects the new argument structure, while the migrated policy completes the authorized action. These checks validate an implementation slice; they do not establish resistance to prompt injection by a real language model. The report covers background, related work, the proposed solution, reproducibility, limitations, and a plan for further evaluation.
\keywords{LLM agents \and indirect prompt injection \and tool contracts \and security regression \and access control}
\end{abstract}
\section{Vấn đề cần giải quyết}
\label{sec:problem}
\subsection{Bối cảnh và động cơ}
Một agent sử dụng mô hình ngôn ngữ lớn (LLM) thường lặp lại ba bước: đọc yêu cầu và dữ liệu, đề xuất lời gọi công cụ, rồi dùng kết quả công cụ để quyết định bước tiếp theo. Khác với hệ thống chỉ sinh văn bản, lời gọi có thể gửi thư, thêm lịch hoặc sửa tệp. Vì vậy, sai lệch trong suy luận có thể dẫn đến thay đổi thực tế ở môi trường.

Prompt injection gián tiếp xuất hiện khi chỉ dẫn của bên thứ ba nằm trong dữ liệu mà agent được phép đọc, thay vì trong yêu cầu trực tiếp của người dùng. Ví dụ, một tài liệu công việc có thể chứa nội dung yêu cầu agent thực hiện hành động ngoài nhiệm vụ. Ranh giới cần bảo vệ là quyền hành động: việc đọc được một chỉ dẫn không đồng nghĩa với việc được phép thực hiện chỉ dẫn đó \cite{greshake2023,debenedetti2024}.

Hệ thống còn phải đối mặt với thay đổi giao diện công cụ. Policy có thể kiểm tra một trường tham số ở phiên bản cũ nhưng không còn diễn giải đúng cấu trúc ở phiên bản mới. Một thay đổi vẫn hợp lệ về schema chưa chắc tương đương về quyền truy cập. Kiểm thử hồi quy cần quan sát cả tác động bị cấm lẫn khả năng hoàn thành tác vụ hợp lệ.

\subsection{Câu hỏi và mục tiêu}
Câu hỏi nghiên cứu là: một phòng thủ đã ngăn tác động có hại ở hợp đồng công cụ $v0$ có duy trì được sự bảo vệ khi chuyển sang $v1$, và việc chuyển đổi policy có khắc phục sai lệch mà vẫn giữ utility hay không? Trong phạm vi học phần, nhóm xây dựng thiết kế đo lường và hiện thực một phần giải pháp để kiểm chứng khả năng triển khai.

Ba đầu ra kỹ thuật của báo cáo là: định nghĩa so sánh theo cặp; xây dựng guard giới hạn quyền theo nhiệm vụ; và kiểm tra độc lập trạng thái sau thực thi. Đề tài cố định có tên tiếng Anh là \emph{Security Regression Testing of Indirect Prompt Injection Defenses under Versioned Agent Tool Contracts}. Báo cáo không khẳng định đã tạo một cơ chế phòng thủ mới hoặc đã chứng minh hiệu quả trên một quần thể agent.

\subsection{Phạm vi và mô hình đe dọa}
Bên tấn công có thể kiểm soát nội dung email, tài liệu hoặc kết quả công cụ mà agent đọc. Họ không được sửa policy tin cậy, dispatcher, mã công cụ hoặc oracle. Môi trường triển khai là sandbox với dữ liệu giả lập. Phạm vi không bao gồm đầu độc mô hình, chiếm quyền mã Python trong tiến trình tin cậy, hoặc bảo đảm luồng thông tin xuyên nhiều agent. Các giả định này giới hạn ý nghĩa của kết quả: guard tại dispatcher không bảo vệ một đường thực thi cố tình bỏ qua dispatcher.

\section{Background}
\subsection{Tool Contracts and Task-Scoped Authorization}
A tool contract specifies operation names, argument schemas, defaults, resource scope, and effect semantics. A rename may preserve meaning, while changing how a resource or destination is represented can alter the inputs to authorization. This report distinguishes the two so that not every interface change is labeled a security regression.

Least privilege grants an agent only the operations needed for its task. A capability may constrain operations, resources, destinations, expiration, and call counts. The implementation in this report covers tools and resources for one task; expiration, quotas, and approval are possible extensions, not verified features.

\subsection{Pre-Execution Enforcement and Post-Execution Measurement}
The guard checks a call before it reaches the executor. ALLOW dispatches the tool call; DENY returns a policy response without executing the tool. If the authorization structure cannot be interpreted, the guard rejects the call, a behavior known as \emph{fail closed}. This can preserve security while reducing utility.

An oracle is an independent scoring function. The security oracle reads pre- and post-execution state to identify forbidden effects; the utility oracle checks whether the task was completed. ALLOW does not prove task success, and DENY alone does not prove that no effect occurred. A tool error is not a security success. Execution errors, policy decisions, and oracle outcomes must therefore be recorded separately.

\subsection{Paired Comparison and Missing Data}
Let $H_{i,v}\in\{0,1\}$ denote the harmful effect for case $i$ under version $v$. A security regression occurs when $H_{i,v0}=0$ and $H_{i,v1}=1$. All four transitions, $00$, $01$, $10$, and $11$, must be retained; reporting only the attack-success rate under $v1$ can hide the distinction between regressions and repairs.
\begin{equation}
\widehat{\Delta}_{\mathrm{sec}}=\frac{1}{n}\sum_{i=1}^{n}(H_{i,v1}-H_{i,v0}).
\end{equation}
Here, $n$ is the number of valid scored pairs. Pairs with missing evidence remain in the planned denominator and are reported as unscorable; they are not replaced with zero. For utility $U$, a regression is $U_{v0}=1$ and $U_{v1}=0$. The two endpoints have different interpretations.

Variants derived from the same source task are not independent observations. A later study must cluster them by source task when estimating uncertainty. This course report does not draw statistical inferences from deterministic checks.

\section{Nghiên cứu liên quan}
AgentDojo cung cấp môi trường có trạng thái để đánh giá agent trước prompt injection và kiểm tra cả security lẫn utility \cite{debenedetti2024}. Nhóm dùng môi trường này làm nguồn tác vụ và công cụ. Việc gọi thành công một công cụ giả lập chỉ xác nhận tích hợp; đánh giá agent cần hành động thực sự do mô hình sinh ra.

InjecAgent nghiên cứu prompt injection gián tiếp trong agent tích hợp công cụ \cite{zhan2024}. Công trình này giúp phân biệt nội dung mang chỉ dẫn tấn công với hành động đầu ra của agent. Báo cáo không gộp một chuỗi lời gọi được viết sẵn với một episode có mô hình đọc carrier và tự quyết định.

CaMeL tổ chức tách biệt dữ liệu không tin cậy và luồng điều khiển, kết hợp kiểm tra capability \cite{debenedetti2025}. Progent đặt trọng tâm vào kiểm soát đặc quyền thông qua policy \cite{shi2025}. Vì các ý tưởng kiểm soát quyền và runtime enforcement đã được nghiên cứu, TaskFileScope trong báo cáo là thành phần tham chiếu để kiểm thử, không phải tuyên bố sáng tạo một nguyên lý bảo vệ mới.

Điểm tập trung của đề tài là đo sự thay đổi của cùng phòng thủ dưới hai hợp đồng công cụ. Đây là một câu hỏi đánh giá có thể kiểm chứng. Tổng hợp trên chỉ là lựa chọn công trình gần, không phải tổng quan hệ thống đầy đủ và chưa đủ để khẳng định tính mới của đóng góp nghiên cứu.

\section{Proposed Solution and Implemented Component}
\subsection{Architecture and Execution Flow}
Trusted inputs consist of the user request, task policy, contract version, and experiment manifest. The agent receives tool data and proposes an action. The dispatcher sends the action through the guard; only an authorized call reaches the executor. Each call records the policy identity, argument digest, decision, dispatch status, and pre/post state hashes. At the end of an episode, an independent oracle scores utility or harmful effects from the retained evidence.

The manifest fixes the planned cases before execution. The capture bridge starts a fresh process to reduce state leakage between episodes and records the runtime version and requests/responses. The validator checks the evidence, and the statistics adapter accepts only valid pairs. Hashes detect changes and support consistency checks; they do not authenticate that a provider generated a response.

\subsection{TaskFileScope and Policy Migration}
The implemented slice uses Workspace task \texttt{user\_task\_29}: read the team-activities file and append activities to it. The trusted policy grants only filename search for the specified file, reading file ID 3, and appending to file ID 3. Other tools are denied. The policy does not derive permissions from document content, attack labels, or model verdicts.

Under $v0$, the read and append tools use a top-level \texttt{file\_id}. Under $v1$, the same tools receive a \texttt{resource} object containing \texttt{space} and \texttt{file\_id}; the space must be \texttt{cloud\_drive}. The adapter validates the structure and maps it to the corresponding native argument. There is no automatic fallback to the old field.

Both configurations belong to the same defense family. F keeps the $v0$ policy unchanged under both contract versions. M uses the same $v0$ policy under $v0$ and changes only the argument decoder to $v1$ under contract $v1$. Tool and resource grants do not change. F and M see the same tool schema within each contract version.

\subsection{Enforcement Rules and Evidence}
The guard rejects ungranted tools, extra arguments, invalid types, incorrect scope, or a different file. DENY is a normal policy response; the native tool is not called. Native tool failures remain execution errors. For a denial with typed state, the pre/post hashes must match. The raw trace records what the agent proposed; the canonical trace normalizes arguments for the upstream scorer.

Filename search retains AgentDojo's substring semantics. This guard is not a system-wide read-confidentiality mechanism. It also does not protect against malicious Python code that can access the callable directly within the process.

\subsection{Relevance to Cybersecurity Governance}
A schema change should be governed like a policy change: assign an owner to review permissions, test before deployment, and retain versioned evidence. If a policy has not been migrated, an explicit denial is clearer than guessing the intended meaning. Organizations must also monitor legitimate tasks that are blocked, because denying every tool call does not meet operational needs.

\section{Triển khai và phương án đánh giá}
\subsection{Môi trường và tổ chức artifact}
Mã nguồn triển khai bằng Python, dùng AgentDojo Workspace v1.2.2 tại commit \texttt{089ed468cf3e}. Môi trường tái lập lịch sử khóa CPython 3.12.13, Pydantic 2.11.4 và PyYAML 6.0.2. Các fixture trạng thái còn phụ thuộc encoding mặc định cp1252 trên Windows với Python UTF-8 mode tắt. Đây là ràng buộc tái lập của dữ liệu hiện có, không phải khuyến nghị dùng cp1252 cho dữ liệu mới.

Guard và adapter migration nằm trong module TaskFileScope; source preflight chạy chuỗi thao tác được cấp quyền qua ToolsExecutor. Manifest và result validator v2 giữ riêng đường kiểm tra đổi tên lịch sử. Amendment TaskFileScope bổ sung schema resource; report acceptance mới đang được củng cố và chưa được xem là kết quả nghiên cứu hoàn chỉnh.

\subsection{Kết quả kiểm chứng thành phần}
Bảng~\ref{tab:utility} ghi kết quả source-conformance đã quan sát. Đầu vào là chuỗi thao tác xác định cho nhiệm vụ người dùng, không phải hành động do LLM thật sinh ra. Mỗi ô bắt đầu từ cùng trạng thái nguồn được reset và utility do scorer upstream tính.
\begin{table}[htbp]
\caption{Utility trong bốn kiểm tra nguồn của TaskFileScope}
\label{tab:utility}
\centering
\begin{tabular}{l@{\hspace{10pt}}l@{\hspace{10pt}}l@{\hspace{10pt}}l}
\toprule
Cấu hình & Contract & Policy decoder & Utility \\
\midrule
F & $v0$ & $v0$ & Thành công \\
F & $v1$ & $v0$ & Không thành công \\
M & $v0$ & $v0$ & Thành công \\
M & $v1$ & $v1$ & Thành công \\
\bottomrule
\end{tabular}
\end{table}

F/$v1$ không hiểu đường dẫn resource mới nên từ chối append. M/$v1$ diễn giải đúng resource và cho phép append vào file được cấp quyền. Kết quả phù hợp với dự đoán về mất utility khi đóng băng policy. Nó không chứng minh có hồi quy an ninh: nhiệm vụ này không cấp quyền xóa ở bất kỳ phiên bản nào, nên lát cắt chưa phân biệt được giả thuyết harmful-effect chính.

Kiểm tra tích hợp với server giả tương thích OpenAI đã chạy lời gọi tìm kiếm và append qua tiến trình capture mới: F/$v1$ giữ nguyên trạng thái khi từ chối; M/$v1$ làm thay đổi trạng thái sau append. Đây là bằng chứng đường chạy và thu thập dữ liệu. Server giả không thể chứng minh mô hình chống được prompt injection.

Bộ regression ở snapshot \texttt{e678ea2} ghi nhận 144 test thành công, không failure hoặc skip. Con số này áp dụng cho snapshot đó; không chứng nhận các thay đổi nghiên cứu đang làm dở sau snapshot. Các test gồm kiểm tra metadata, trace, reset, thiếu bằng chứng và request retry. Báo cáo không dùng số test làm thước đo hiệu quả bảo mật.

\subsection{Tái lập và tiêu chí nghiệm thu}
Trong môi trường pinned của repository, kiểm tra nguồn chạy bằng module source preflight của TaskFileScope; kiểm tra tích hợp nằm trong bộ test amendment cùng tên. Hướng dẫn đi kèm bộ LaTeX ghi đường dẫn và lệnh PowerShell đầy đủ. Biên dịch báo cáo không cần môi trường AgentDojo hoặc khóa API.

Nghiệm thu thành phần yêu cầu: lời gọi ngoài quyền không tới native executor; denial không đổi typed state; policy F không tự đổi theo contract; M chỉ thay decoder đã khai báo; utility được tính độc lập. Thất bại provider, lỗi reset hoặc thiếu trace phải được giữ trong mẫu dự kiến, không được xóa để làm đẹp kết quả.

\subsection{Thiết kế thực nghiệm tiếp theo}
Giai đoạn nghiên cứu cần bổ sung một treatment có khả năng bộc lộ thay đổi quyền tạo tác động có hại, không làm yếu F chỉ để tạo kết quả mong muốn. Với mỗi case, giữ cùng model snapshot, task, carrier và reset; chạy cả $v0$ và $v1$ cho F/M. Tách riêng attack endpoint và default-control utility; lưu đủ bốn chuyển tiếp và nguyên nhân không chấm được.

Khi có đủ dữ liệu, báo cáo chênh lệch rủi ro theo cặp, khoảng tin cậy bootstrap theo nhiệm vụ nguồn và McNemar trong điều kiện độc lập thích hợp. Kế hoạch nghiên cứu dự kiến ít nhất ba model snapshots, kiểm tra carrier giữ lại và tái lập độc lập. Chưa có bộ dữ liệu đáp ứng các cổng đó trong kết quả của báo cáo này.

\section{Hạn chế và kết luận}
Phần hiện thực đã cho thấy có thể đặt một guard quyền theo nhiệm vụ trước công cụ, chuyển đổi cấu trúc resource một cách tường minh và kiểm tra utility bằng scorer độc lập. Đây là phần giải pháp cụ thể phục vụ mục tiêu học phần; chưa phải nghiên cứu thực nghiệm hoàn chỉnh về khả năng chống prompt injection.

Các hạn chế chính gồm một task nguồn, một kiểu migration do nhóm thiết kế, chuỗi thao tác xác định và provider giả. Thiết kế chưa kiểm tra paraphrase, partial leakage, công cụ ngoài sandbox hoặc nhiều model. Biến thiên timestamp của thao tác append cũng khiến hash trạng thái sau thực thi có thể khác, dù utility giống nhau. Không được suy rộng số đếm kiểm tra sang tỷ lệ an toàn của hệ thống thực tế.

Bước tiếp theo về kỹ thuật là hoàn thiện ràng buộc plan--record--oracle và kiểm chứng phản biện độc lập. Sau đó mới thu thập dữ liệu từ mô hình thật theo thiết kế cố định. Kết luận hiện tại là tính khả thi của phương pháp kiểm thử và một lát cắt triển khai đã được kiểm tra; kết luận hiệu quả bảo mật vẫn cần bằng chứng bổ sung.

\clearpage
\renewcommand{\refname}{References}
\bibliographystyle{splncs04}
\bibliography{references}
\end{document}
